% TOP_PROBLEM: {"id": 3104, "title": "Edge Reconstruction Conjecture", "category_id": 3, "category": {"id": 3, "name": "graph_theory", "display_name": "Graph Theory", "description": "Problems involving graphs, networks, and their properties.", "slug": "graph-theory", "order_index": 3, "created_at": "2026-07-31T15:26:25.671Z"}, "status": "open", "problem_number": "OPG-804", "set_id": 12}
% FIRST_POSTED: 2026-10-01
% ENTRY_KIND: statement_only
% UnsolvedMath ID 3104; source catalogue code OPG-804
% !TeX program = lualatex
% Definition and sources only; this file is not an LLM solution attempt.
\documentclass[11pt,a4paper]{article}
\usepackage[margin=25mm]{geometry}
\usepackage{fontspec}
\setmainfont{lmroman10-regular.otf}[BoldFont=lmroman10-bold.otf,
 ItalicFont=lmroman10-italic.otf,BoldItalicFont=lmroman10-bolditalic.otf]
\usepackage{amsmath,amssymb,mathtools}
\usepackage{unicode-math}
\setmathfont{latinmodern-math.otf}
\AtBeginDocument{\let\setminus\smallsetminus}
\usepackage{xurl,hyperref}
\hypersetup{colorlinks=true,urlcolor=blue}
\setcounter{secnumdepth}{0}
\setlength{\parindent}{0pt}
\setlength{\parskip}{5pt}
\setlength{\emergencystretch}{3em}
\begin{document}
\section*{Edge Reconstruction Conjecture}\label{3104}
{\small UnsolvedMath ID 3104 · OPG-804 · Graph Theory · OpenGarden}
\subsection{Definitions and mathematical statement}
Let $G=(V,E)$ be a finite simple undirected graph with $m=|E|$.
For an edge $e\in E$, define $G-e=(V,E\setminus\{e\})$.
Only that edge is deleted: every vertex is retained, including endpoints
that become isolated and vertices already isolated. The edge deck is the
multiset of unlabelled isomorphism classes
\[
 \mathcal D_E(G)=\{\!\{[G-e]:e\in E\}\!\}.
\]
Each original edge contributes one card, even if several deletions produce
isomorphic graphs. The cards carry no common vertex labels and do not
identify which edge was removed.

The conjecture asserts that whenever finite simple graphs $G,H$ each
have at least four edges,
\[
 \mathcal D_E(G)=\mathcal D_E(H)\quad\Longrightarrow\quad G\cong H.
\]
Here an isomorphism is a bijection of vertex sets preserving adjacency.
The hypothesis gives exactly the whole multiset of cards, not only the
set of distinct isomorphism types. Its size records the original number
of edges, and the vertex count of every card records the original order.

There is no connectedness assumption, and isolated vertices remain part
of the object to reconstruct. The lower threshold refers to four edges,
not four vertices. The question concerns reconstruction from edge deletion;
removing an endpoint as well would produce a different deck.
An affirmative answer needs only uniqueness up to isomorphism, whereas
a negative answer requires two genuinely nonisomorphic graphs satisfying
the complete multiset equality.
\subsection{Short English statement}
Is every finite simple graph with at least four edges determined up to isomorphism by all its edge-deleted cards, counted with multiplicity?
\subsection{Sources}
\begin{itemize}
\item[S1] ULAM.AI and the UnsolvedMath Contributors, supplied problems.json, record 3104 (OPG-804), OpenGarden collection. Catalogue identity and source transcription; CC BY 4.0.
\url{https://huggingface.co/datasets/ulamai/UnsolvedMath}.
\item[S2] Open Problem Garden, Edge Reconstruction Conjecture. Original problem and discussion; the mathematical formulation below is expanded from the supplied source record.
\url{http://www.openproblemgarden.org/op/edge_reconstruction_conjecture}.
\end{itemize}
\end{document}
